\documentclass[11pt,a4paper]{scrartcl}


\usepackage[utf8]{inputenc}
\usepackage[T1]{fontenc}
\newcommand{\ats}[1]{\par\noindent\textit{Atsakymas:} #1}
\usepackage{cancel}
\usepackage{tikz}
\usepackage{lmodern} 
\usepackage{amsmath,amssymb,amsthm}
\usepackage{graphicx}
\usepackage[dvipsnames,svgnames]{xcolor}
\usepackage{hyperref}
\usepackage{mathrsfs}
\usepackage[shortlabels]{enumitem}
\usepackage[obeyFinal,textsize=scriptsize,shadow,loadshadowlibrary]{todonotes}
\usepackage{mathtools}
\usepackage{microtype}

%%% Paragraph layout
\setlength{\parindent}{0pt}
\setlength{\parskip}{0.6\baselineskip}

%%% Colored sections
\renewcommand*{\sectionformat}%
  {\color{purple}\S\thesection\enskip}
\renewcommand*{\subsectionformat}%
  {\color{purple}\S\thesubsection\enskip}
\renewcommand*{\subsubsectionformat}%
  {\color{purple}\S\thesubsubsection\enskip}
\addtokomafont{paragraph}{\color{orange!35!black}\P\ }
\KOMAoptions{numbers=noenddot}

%%% Title
\addtokomafont{subtitle}{\Large}
\setkomafont{author}{\Large\scshape}
\setkomafont{date}{\Large\normalsize}
% Neigiamas vertikalus tarpas, kuris pakelia dokumento antraštę į viršų. 
% Galite keisti -2.5cm į -3cm (aukščiau) arba -1.5cm (žemiau).
\titlehead{\vspace*{-7.5cm}} 

% PRIDĖTA: Automatiškai perrašome datą į mokyklos pavadinimą
\AtBeginDocument{%
  \date{\normalsize Vilniaus licėjus}
}

%%% Page setup
\usepackage[headsepline]{scrlayer-scrpage}
\renewcommand{\headfont}{}
\addtolength{\textheight}{3.14cm}
\setlength{\footskip}{0.5in}
\setlength{\headsep}{10pt}
% Puslapio antraštė turi rodyti tik konkrečius dokumento duomenis.
% Nustatome juos aiškiai, kad neatsirastų "\@author" / "\@title" arba senų reikšmių.
\makeatletter
\AtBeginDocument{%
  \ihead{\footnotesize\textbf{Andrius Berniukevičius}}%
  \ohead{\footnotesize\textbf{\@title}}%
}
\makeatother
\automark{section}
\chead{}
\cfoot{\pagemark}

%%% Macros
\providecommand{\ol}{\overline}
\providecommand{\ul}{\underline}
\providecommand{\wt}{\widetilde}
\providecommand{\wh}{\widehat}
\providecommand{\eps}{\varepsilon}
\providecommand{\half}{\frac{1}{2}}
\providecommand{\inv}{^{-1}}
\newcommand{\dang}{\measuredangle} %% Directed angle
\newcommand{\vocab}[1]{\textbf{\color{blue}\sffamily #1}}
\providecommand{\CC}{\mathbb C}
\providecommand{\FF}{\mathbb F}
\providecommand{\NN}{\mathbb N}
\providecommand{\QQ}{\mathbb Q}
\providecommand{\RR}{\mathbb R}
\providecommand{\ZZ}{\mathbb Z}
\providecommand{\ts}{\textsuperscript}
\providecommand{\dg}{^\circ}
\providecommand{\ii}{\item}
\providecommand{\defeq}{\coloneqq}
\DeclareMathOperator*{\lcm}{lcm}
\DeclareMathOperator*{\argmin}{arg min}
\DeclareMathOperator*{\argmax}{arg max}
\providecommand{\hrulebar}{\par
\hspace{\fill}\rule{0.95\linewidth}{.7pt}\hspace{\fill}
\par\nointerlineskip \vspace{\baselineskip}}

%%% Optional colored theorems
\usepackage{thmtools}
\usepackage[framemethod=TikZ]{mdframed}
\mdfdefinestyle{mdbluebox}{%
  roundcorner=10pt,
  linewidth=1pt,
  skipabove=12pt,
  innerbottommargin=9pt,
  skipbelow=2pt,
  linecolor=blue,
  nobreak=true,
  backgroundcolor=TealBlue!5,
}
\declaretheoremstyle[
  headfont=\sffamily\bfseries\color{MidnightBlue},
  mdframed={style=mdbluebox},
  headpunct={\\[3pt]},
  postheadspace={0pt}
]{thmbluebox}
\mdfdefinestyle{mdredbox}{%
  linewidth=0.5pt,
  skipabove=12pt,
  frametitleaboveskip=5pt,
  frametitlebelowskip=0pt,
  skipbelow=2pt,
  frametitlefont=\bfseries,
  innertopmargin=4pt,
  innerbottommargin=8pt,
  nobreak=true,
  backgroundcolor=Salmon!5,
  linecolor=RawSienna,
}
\declaretheoremstyle[
  headfont=\bfseries\color{RawSienna},
  mdframed={style=mdredbox},
  headpunct={\\[3pt]},
  postheadspace={0pt},
]{thmredbox}
\mdfdefinestyle{mdgreenbox}{%
  skipabove=8pt,
  linewidth=2pt,
  rightline=false,
  leftline=true,
  topline=false,
  bottomline=false,
  linecolor=ForestGreen,
  backgroundcolor=ForestGreen!5,
}
\declaretheoremstyle[
  headfont=\bfseries\sffamily\color{ForestGreen!70!black},
  bodyfont=\normalfont,
  spaceabove=2pt,
  spacebelow=1pt,
  mdframed={style=mdgreenbox},
  headpunct={ --- },
]{thmgreenbox}
\mdfdefinestyle{mdblackbox}{%
  skipabove=8pt,
  linewidth=3pt,
  rightline=false,
  leftline=true,
  topline=false,
  bottomline=false,
  linecolor=black,
  backgroundcolor=RedViolet!5!gray!5,
}
\declaretheoremstyle[
  headfont=\bfseries,
  bodyfont=\normalfont\small,
  spaceabove=0pt,
  spacebelow=0pt,
  mdframed={style=mdblackbox}
]{thmblackbox}

%% Aplinkos su rėmeliais (thmtools)
\declaretheorem[style=thmbluebox,name=Teorema]{theorem}
\declaretheorem[style=thmbluebox,name=Lema]{lemma}
\declaretheorem[style=thmbluebox,name=Savybė]{property}
\declaretheorem[style=thmbluebox,name=Teiginys]{proposition}
\declaretheorem[style=thmbluebox,name=Išvada]{corollary}
\declaretheorem[style=thmbluebox,name=Teorema,numbered=no]{theorem*}
\declaretheorem[style=thmbluebox,name=Lema,numbered=no]{lemma*}
\declaretheorem[style=thmbluebox,name=Teiginys,numbered=no]{proposition*}
\declaretheorem[style=thmbluebox,name=Išvada,numbered=no]{corollary*}
\declaretheorem[style=thmgreenbox,name=Algoritmas]{algorithm}
\declaretheorem[style=thmgreenbox,name=Algoritmas,numbered=no]{algorithm*}
\declaretheorem[style=thmgreenbox,name=Teiginys]{claim}
\declaretheorem[style=thmgreenbox,name=Teiginys,numbered=no]{claim*}
\declaretheorem[style=thmredbox,name=Pavyzdys]{example}
\declaretheorem[style=thmredbox,name=Pavyzdys,numbered=no]{example*}
\declaretheorem[style=thmblackbox,name=Pastaba]{remark}
\declaretheorem[style=thmblackbox,name=Pastaba,numbered=no]{remark*}

%% Standartinės amsthm aplinkos (Apibrėžimai ir užduotys)
\theoremstyle{definition}
\ifcsname conjecture\endcsname\else\newtheorem{conjecture}{Hipotezė}\fi
\ifcsname definition\endcsname\else\newtheorem{definition}{Apibrėžimas}\fi
\ifcsname fact\endcsname\else\newtheorem{fact}{Faktas}\fi
\ifcsname answer\endcsname\else\newtheorem{answer}{Atsakymas}\fi
\ifcsname ques\endcsname\else\newtheorem{ques}{Klausimas}\fi
\ifcsname exercise\endcsname\else\newtheorem{exercise}{Užduotis}\fi
\ifcsname problem\endcsname\else\newtheorem{problem}{Uždavinys}[section]\fi
\ifcsname conjecture*\endcsname\else\newtheorem*{conjecture*}{Hipotezė}\fi
\ifcsname definition*\endcsname\else\newtheorem*{definition*}{Apibrėžimas}\fi
\ifcsname fact*\endcsname\else\newtheorem*{fact*}{Faktas}\fi
\ifcsname answer*\endcsname\else\newtheorem*{answer*}{Atsakymas}\fi
\ifcsname ques*\endcsname\else\newtheorem*{ques*}{Klausimas}\fi
\ifcsname exercise*\endcsname\else\newtheorem*{exercise*}{Užduotis}\fi
\ifcsname problem*\endcsname\else\newtheorem*{problem*}{Uždavinys}\fi

\newcounter{answerssection}
\newcommand{\answerssection}{%
  \setcounter{answerssection}{\value{section}}%
  \section{Atsakymai}%
  \setcounter{problem}{0}%
  \renewcommand{\theproblem}{\arabic{answerssection}.\arabic{problem}}%
}

%% GALUTINIS SPRENDIMAS PROOF APLINKAI
%% --- PRADŽIA: Įrodymo aplinkos sutvarkymas ---
\makeatletter

% 1. Užtikriname, kad žodis būtų "Įrodymas", net jei "babel" paketas bando jį perrašyti
\AtBeginDocument{%
  \renewcommand{\proofname}{Įrodymas}%
  \@ifpackageloaded{babel}{%
    \addto\captionslithuanian{\renewcommand{\proofname}{Įrodymas}}%
  }{}%
}

% 2. Priverstinai pašaliname scrartcl klasės bold šriftą ir paliekame TIK italic
% 3. Sujungiame su tuščiu kvadratėliu dešiniajame kampe.
\renewcommand{\qedsymbol}{\ensuremath{\Box}}
\renewenvironment{proof}[1][\proofname]{\par
  \pushQED{\qed}%
  \normalfont \topsep6\p@\@plus6\p@\relax
  \trivlist
  \item[\hskip\labelsep
        \normalfont\itshape % Priverčia naudoti tik pasvirą šriftą, kaip nuotraukoje
    #1\@addpunct{.}]\ignorespaces
}{%
  \popQED\endtrivlist\@endpefalse
}
\newenvironment{solution}{\par
  \normalfont \topsep6\p@\@plus6\p@\relax
  \trivlist
  \item[\hskip\labelsep
        \normalfont\itshape
    \textit{Sprendimas}\@addpunct{.}]\ignorespaces
}{%
  \endtrivlist\@endpefalse
}
\newcommand{\ans}[1]{\par\noindent\textit{Atsakymas:} #1}
\makeatother


\title{Dalumas, lyginumas ir DBD}
\author{Andrius Berniukevičius}

\newenvironment{solution}
    {\begin{list}{}{\setlength{\leftmargin}{10pt}\setlength{\rightmargin}{10pt}}\item[]}
    {\end{list}}

\begin{document}

\maketitle

% ============================================================
\section{Dalumo žymėjimai ir kalba}
% ============================================================

\subsection{Ką reiškia $b\mid a$?}

Tegu $a,b\in\mathbb Z$ ir $b\neq 0$. Žymėjimas
\[
b\mid a
\]
reiškia, kad \textbf{$b$ dalija $a$}, t.~y. egzistuoja toks sveikasis skaičius
$k$, kad
\[
a=bk.
\]
Pavyzdžiui,
\[
4\mid 20,
\]
nes
\[
20=4\cdot 5.
\]
Tačiau
\[
4\nmid 22,
\]
nes nėra sveikojo skaičiaus $k$, kuriam
\[
22=4k.
\]
Svarbu nepainioti:
\[
b\mid a
\]
reiškia, kad \textbf{$b$ yra daliklis, o $a$ --- dalinys}. 
Jeigu
\[
a=bk,
\]
tai taip pat sakome, kad $a$ yra $b$ kartotinis.

\subsection{Naudingos dalumo savybės}

Tegu $a,b,c,m,n\in\mathbb Z$.

\begin{enumerate}
    \renewcommand{\labelenumi}{\textbf{\theenumi.}}

    \item Jei $a\mid b$, tai $a\mid bc$.

    \textbf{Įrodymas.} Kadangi $a\mid b$, egzistuoja $k\in\mathbb Z$, kad
    $b=ak$. Tada
    \[
    bc=(ak)c=a(kc).
    \]
    Kadangi $kc\in\mathbb Z$, tai $a\mid bc$.
    \hfill $\square$

    \item Jei $a\mid b$ ir $a\mid c$, tai
    \[
    a\mid (b+c)
    \qquad\text{ir}\qquad
    a\mid (b-c).
    \]

    \textbf{Įrodymas.} Iš $a\mid b$ gauname $b=am$, o iš $a\mid c$ ---
    $c=an$ kai $m,n\in\mathbb Z$. Tada
    \[
    b+c=am+an=a(m+n),
    \]
    \[
    b-c=am-an=a(m-n).
    \]
    Kadangi $m\pm n\in\mathbb Z$, gauname $a\mid (b+c)$ ir $a\mid (b-c)$.
    \hfill $\square$

    \item Dar bendriau: jei $a\mid b$ ir $a\mid c$, tai bet kuriems
    $m,n\in\mathbb Z$
    \[
    a\mid (mb+nc).
    \]

    \textbf{Įrodymas.} Kaip ir aukščiau, $b=am$ ir $c=an$ kai $m,n\in\mathbb Z$.
    Tuomet
    \[
    mb+nc=(am)u+(an)v=a(mu+nv),
    \]
    kur $u,v\in\mathbb Z$ yra atitinkami koeficientai. Kadangi $mu+nv\in\mathbb Z$,
    gauname \[a\mid (mb+nc)\].
    \hfill $\square$

    \item Jei $a\mid b$ ir $b\mid c$, tai $a\mid c$.

    \textbf{Įrodymas.} Iš $a\mid b$ gauname $b=ak$, o iš $b\mid c$ ---
    $c=bl$ su $k,l\in\mathbb Z$. Tada
    \[
    c=bl=(ak)l=a(kl).
    \]
    Kadangi $kl\in\mathbb Z$, tai $a\mid c$.
    \hfill $\square$

    \item Jei $a\mid b$ ir $b\mid a$, tai
    \[
    |a|=|b|.
    \]

    \textbf{Įrodymas.} Iš $a\mid b$ gauname $b=ak$ su $k\in\mathbb Z$.
    Iš $b\mid a$ gauname $a=bl$ su $l\in\mathbb Z$. Tuomet
    \[
    a=bl=(ak)l=a(kl).
    \]
    Jei $a\neq 0$, tai $kl=1$, todėl $k=\pm 1$, o tai reiškia $|b|=|a|$.
    Jei $a=0$, tada $b=0$, ir vėl $|a|=|b|=0$.
    \hfill $\square$
\end{enumerate}

Trečioji savybė ypač svarbi olimpiadose. Ji sako, kad jeigu skaičius dalija du
skaičius, tai jis dalija ir bet kurią jų \textbf{sveikųjų koeficientų tiesinę kombinaciją}.

\begin{example}

Žinoma, kad $7\mid a$ ir $7\mid b$. Įrodykite, kad
\[
7\mid (3a-5b).
\]

\end{example}

\begin{solution}
\textbf{Sprendimas}

Kadangi $7\mid a$, egzistuoja $m\in\mathbb Z$, kuriam
\[
a=7m.
\]

Kadangi $7\mid b$, egzistuoja $n\in\mathbb Z$, kuriam
\[
b=7n.
\]

Tuomet
\[
3a-5b=3\cdot 7m-5\cdot 7n=7(3m-5n).
\]

Kadangi $3m-5n\in\mathbb Z$, gauname
\[
7\mid(3a-5b).
\]

\hfill $\square$

\end{solution}

\subsection{Išvada olimpiadininkui}

Jeigu sąlygoje matai kelis dalumo teiginius, verta paklausti:

\begin{center}
\textbf{„Ar sudėjęs, atėmęs arba padauginęs šiuos skaičius iš tinkamų
sveikųjų skaičių galiu gauti ką nors paprastesnio?“}
\end{center}

% ============================================================
\section{Dalumas ir liekana be modulių kalbos}
% ============================================================

Jeigu norime išsiaiškinti, kada išraiška $A(n)$ dalijasi iš $B(n)$, dažnai
naudinga perrašyti
\[
A(n)=B(n)\cdot Q(n)+r,
\]
kur $r$ yra mažas skaičius.

Tuomet
\[
B(n)\mid A(n)
\]
galima pakeisti daug paprastesne sąlyga
\[
B(n)\mid r.
\]

\begin{example}

Raskite visus sveikuosius skaičius $n$, kuriems
\[
n+2\mid n^2+5n+8.
\]

\end{example}

\begin{solution}
\textbf{Sprendimas}

Pertvarkome:
\[
n^2+5n+8=(n+2)(n+3)+2.
\]

Jeigu
\[
n+2\mid n^2+5n+8,
\]
tai $n+2$ turi dalyti ir skirtumą
\[
\bigl(n^2+5n+8\bigr)-(n+2)(n+3)=2.
\]

Taigi
\[
n+2\mid 2.
\]

Sveikieji skaičiaus $2$ dalikliai yra
\[
\pm1,\ \pm2.
\]

Todėl
\[
n+2\in\{-2,-1,1,2\},
\]
taigi
\[
n\in\{-4,-3,-1,0\}.
\]

Visos keturios reikšmės tinka.

\textbf{Atsakymas:}
\[
\boxed{n\in\{-4,-3,-1,0\}}.
\]

\hfill $\square$

\end{solution}

% ============================================================
\section{Lyginiai ir nelyginiai skaičiai}
% ============================================================

Sveikasis skaičius yra \textbf{lyginis}, jeigu jis dalijasi iš $2$. Taigi lyginį skaičių galime užrašyti forma
\[
2k,
\qquad k\in\mathbb Z.
\]
Sveikasis skaičius yra \textbf{nelyginis}, jeigu jis nesidalija iš $2$.
Kiekvieną nelyginį skaičių galima užrašyti forma
\[
2k+1,
\qquad k\in\mathbb Z.
\]
Naudinga mokėti ne tik prisiminti, bet ir greitai įrodyti šias savybes:

\begin{enumerate}
    \renewcommand{\labelenumi}{\textbf{\theenumi.}}

    \item lyginis $+$ lyginis $=$ lyginis;
    \item nelyginis $+$ nelyginis $=$ lyginis;
    \item lyginis $+$ nelyginis $=$ nelyginis;
    \item lyginis $\cdot$ bet koks sveikasis $=$ lyginis;
    \item nelyginis $\cdot$ nelyginis $=$ nelyginis.
\end{enumerate}

\begin{example}

Įrodykite, kad kiekvienam sveikajam skaičiui $n$ skaičius
\[
n(n+1)
\]
yra lyginis.

\end{example}

\begin{solution}
\textbf{Sprendimas}

Skaičiai $n$ ir $n+1$ yra du iš eilės einantys sveikieji skaičiai.
Vienas iš jų būtinai yra lyginis. Todėl bent vienas iš sandaugos
\[
n(n+1)
\]
daugiklių dalijasi iš $2$.

Vadinasi,
\[
2\mid n(n+1).
\]

\end{solution}

\begin{example}

Įrodykite: jeigu $n^2$ yra lyginis, tai ir $n$ yra lyginis.

\end{example}

\begin{solution}
\textbf{Sprendimas}

Įrodysime kontrapoziciją. Tarkime, kad $n$ yra nelyginis. Tuomet
\[
n=2k+1
\]
su tam tikru $k\in\mathbb Z$.

Tada
\[
n^2=(2k+1)^2=4k^2+4k+1
=2(2k^2+2k)+1.
\]

Vadinasi, $n^2$ yra nelyginis.

Taigi, jeigu $n^2$ yra lyginis, $n$ negali būti nelyginis, todėl $n$ yra
lyginis.

\end{solution}

\subsection{Išvada olimpiadininkui}

Kai uždavinyje yra sveikieji skaičiai ir klausiama, ar kažkas
\textbf{įmanoma}, pirmiausia verta patikrinti lyginumą.

Kartais visas ilgas uždavinys subyra į vieną klausimą:

\begin{center}
\textbf{„Kas čia yra lyginis, o kas nelyginis?“}
\end{center}

% ============================================================
\section{Bendrieji dalikliai ir DBD}
% ============================================================

Skaičius $d$ vadinamas skaičių $a$ ir $b$ \textbf{bendruoju dalikliu}, jeigu
\[
d\mid a
\qquad\text{ir}\qquad
d\mid b.
\]
Didžiausias teigiamas bendrasis daliklis vadinamas
\textbf{didžiausiuoju bendruoju dalikliu}. Žymėsime
\[
\operatorname{DBD}(a,b).
\]
Pavyzdžiui,
\[
\operatorname{DBD}(18,30)=6.
\]
Jeigu
\[
\operatorname{DBD}(a,b)=1,
\]
sakome, kad $a$ ir $b$ yra \textbf{tarpusavyje pirminiai}. Pavyzdžiui,
\[
\operatorname{DBD}(8,15)=1.
\]

Tai nereiškia, kad $8$ arba $15$ turi būti pirminiai skaičiai.
Tai reiškia tik tai, kad jie neturi bendro daliklio, didesnio už $1$.

% ============================================================
\section{Svarbiausia DBD idėja: skirtumas nekeičia bendrųjų daliklių}
% ============================================================

Jeigu skaičius $d$ dalija $a$ ir $b$, tai jis dalija ir
\[
a-b.
\]
Todėl skaičių $a$ ir $b$ bendrieji dalikliai yra tokie patys kaip skaičių
$a-b$ ir $b$ bendrieji dalikliai. Taigi
\[
\boxed{\operatorname{DBD}(a,b)=\operatorname{DBD}(a-b,b)}.
\]
Dar bendriau, bet kuriam $q\in\mathbb Z$,
\[
\boxed{\operatorname{DBD}(a,b)=\operatorname{DBD}(a-qb,b)}.
\]
\textbf{Įrodymas.} Tegul $q\in\mathbb Z$.
Tarkime, kad $c$ yra bet koks sveikasis skaičius. Tada
\[
 c\mid a \text{ ir } c\mid b
\Longleftrightarrow
 c\mid (a-qb) \text{ ir } c\mid b.
\]
\begin{itemize}
    \item jei $c\mid a$ ir $c\mid b$, tai $c\mid qb$, o tada
    \[
    c\mid (a-qb),
    \]
    nes $a-qb=a+(-q)b$.

    \item jei $c\mid (a-qb)$ ir $c\mid b$, tai $c\mid qb$, o tada
    \[
    a=(a-qb)+qb,
    \]
    taigi $c\mid a$.
\end{itemize}
Taigi skaičių $a$ ir $b$ bendrųjų daliklių aibė sutampa su skaičių $a-qb$ ir $b$ bendrųjų daliklių aibe. Todėl jų didžiausias bendrasis daliklis yra vienodas:
\[
\operatorname{DBD}(a,b)=\operatorname{DBD}(a-qb,b).
\]
\hfill $\square$

Ši savybė yra Euklido algoritmo pagrindas.

\begin{example}

Apskaičiuokite
\[
\operatorname{DBD}(252,198).
\]

\end{example}

\begin{solution}
\textbf{Sprendimas}

Naudojame Euklido algoritmą:
\[
252=198\cdot1+54,
\]
todėl
\[
\operatorname{DBD}(252,198)=\operatorname{DBD}(198,54).
\]

Toliau
\[
198=54\cdot3+36,
\]
todėl
\[
\operatorname{DBD}(198,54)=\operatorname{DBD}(54,36).
\]

Galiausiai
\[
54=36\cdot1+18,
\]
\[
36=18\cdot2.
\]

Taigi
\[
\boxed{\operatorname{DBD}(252,198)=18}.
\]

\end{solution}

% ============================================================
\section{Kodėl du iš eilės einantys skaičiai yra tarpusavyje pirminiai?}
% ============================================================

\begin{example}

Įrodykite, kad bet kuriam sveikajam $n$ skaičiai $n$ ir $n+1$ yra
tarpusavyje pirminiai.

\end{example}

\begin{solution}
\textbf{Sprendimas}

Tegu $d$ yra bet kuris bendras skaičių $n$ ir $n+1$ daliklis.
Tuomet
\[
d\mid n
\qquad\text{ir}\qquad
d\mid(n+1).
\]

Vadinasi, $d$ dalija ir jų skirtumą:
\[
d\mid\bigl((n+1)-n\bigr).
\]

Taigi
\[
d\mid1.
\]

Vienintelis teigiamas skaičiaus $1$ daliklis yra $1$.
Todėl
\[
\boxed{\operatorname{DBD}(n,n+1)=1}.
\]

\end{solution}

\subsection{Labai svarbus mąstymo šablonas}

Jeigu reikia rasti arba apriboti
\[
\operatorname{DBD}(A,B),
\]
bandyk sudaryti paprastesnę tiesinę kombinaciją
\[
mA+nB.
\]
Bet kuris bendras $A$ ir $B$ daliklis dalys ir šią kombinaciją. Jei pavyksta gauti, pavyzdžiui,
\[
mA+nB=5,
\]
tai iš karto žinome, kad
\[
\operatorname{DBD}(A,B)\mid5.
\]
Vadinasi, DBD gali būti tik $1$ arba $5$.

\begin{example}

Įrodykite, kad bet kuriam sveikajam $n$
\[
\operatorname{DBD}(2n+1,n+3)
\]
gali būti tik $1$ arba $5$.

\end{example}

\begin{solution}
\textbf{Sprendimas}

Tegu
\[
d=\operatorname{DBD}(2n+1,n+3).
\]

Tuomet
\[
d\mid(2n+1)
\qquad\text{ir}\qquad
d\mid(n+3).
\]

Todėl $d$ dalija ir tiesinę kombinaciją
\[
2(n+3)-(2n+1)=5.
\]

Taigi
\[
d\mid5.
\]

Kadangi $d$ yra teigiamas, gauname
\[
\boxed{d\in\{1,5\}}.
\]

\end{solution}

% ============================================================
\section{Dar viena naudinga savybė}
% ============================================================

Jeigu
\[
\operatorname{DBD}(a,b)=1
\]
ir
\[
a\mid bc,
\]
tai
\[
a\mid c.
\]
Šis teiginys vadinamas \textbf{Euklido lema}. Intuityviai: jeigu $a$ neturi bendrų pirminių daliklių su $b$, tai visi
$a$ pirminiai dalikliai, esantys sandaugoje $bc$, turi ateiti iš $c$. Pavyzdžiui, kadangi
\[
\operatorname{DBD}(8,15)=1
\]
ir
\[
8\mid 15c,
\]
tai būtinai
\[
8\mid c.
\]
Šią savybę naudosime vėlesniuose skaičių teorijos uždaviniuose.

% ============================================================
\section{Ką daryti pamačius dalumo uždavinį?}
% ============================================================

Prieš pradėdamas skaičiuoti, paklausk:

\begin{enumerate}
    \renewcommand{\labelenumi}{\textbf{\theenumi.}}

    \item Ar galiu užrašyti dalumo sąlygą forma
    \[
    a=bk?
    \]

    \item Ar galiu sudėti arba atimti du dalumo teiginius?

    \item Ar galiu sudaryti naudingą tiesinę kombinaciją?

    \item Ar galima didelę išraišką perrašyti forma
    \[
    \text{daliklis}\cdot\text{kažkas}+\text{mažas skaičius}?
    \]

    \item Ar svarbus lyginumas?

    \item Ar verta nagrinėti DBD ir iš skaičių skirtumo gauti mažesnį skaičių?

    \item Ar du skaičiai negali būti tarpusavyje pirminiai?
\end{enumerate}

% ============================================================
\newpage
\section{Uždaviniai savarankiškam sprendimui}
% ============================================================

\subsection{A lygis. Dalumo kalba}

\begin{enumerate}
    \renewcommand{\labelenumi}{\textbf{\theenumi.}}

    \item Užrašykite lygybe su sveikuoju skaičiumi $k$, ką reiškia teiginys
    \[
    7\mid a.
    \]

    \item Žinoma, kad $5\mid a$ ir $5\mid b$. Įrodykite, kad
    \[
    5\mid(2a+3b).
    \]

    \item Žinoma, kad $9\mid a$ ir $9\mid b$. Įrodykite, kad
    \[
    9\mid(4a-7b).
    \]

    \item Įrodykite: jei $6\mid a$, tai
    \[
    3\mid a
    \qquad\text{ir}\qquad
    2\mid a.
    \]

    \item Įrodykite: jei $a\mid b$ ir $b\mid c$, tai $a\mid c$.

\end{enumerate}

\subsection{B lygis. Lyginumas}

\begin{enumerate}
    \renewcommand{\labelenumi}{\textbf{\theenumi.}}
    \setcounter{enumi}{5}

    \item Įrodykite, kad dviejų iš eilės einančių sveikųjų skaičių suma yra
    nelyginė.

    \item Įrodykite, kad trijų iš eilės einančių sveikųjų skaičių sandauga
    yra lyginė.

    \item Įrodykite: jeigu $a+b$ yra nelyginis, tai $a$ ir $b$ yra skirtingo
    lyginumo.

    \item Įrodykite: jeigu $ab$ yra nelyginis, tai ir $a$, ir $b$ yra
    nelyginiai.

    \item Įrodykite, kad skaičius
    \[
    n^2+n
    \]
    yra lyginis su kiekvienu sveikuoju $n$.

\end{enumerate}

\subsection{C lygis. DBD ir Euklido algoritmas}

\begin{enumerate}
    \renewcommand{\labelenumi}{\textbf{\theenumi.}}
    \setcounter{enumi}{10}

    \item Euklido algoritmu apskaičiuokite
    \[
    \operatorname{DBD}(414,662).
    \]

    \item Apskaičiuokite
    \[
    \operatorname{DBD}(2026,2025).
    \]
    Paaiškinkite, kodėl atsakymą galima gauti beveik be skaičiavimų.

    \item Įrodykite, kad
    \[
    \operatorname{DBD}(n,n+2)
    \]
    gali būti tik $1$ arba $2$.

    \item Įrodykite, kad
    \[
    \operatorname{DBD}(3n+1,n)=1
    \]
    su kiekvienu sveikuoju $n$.

    \item Įrodykite, kad
    \[
    \operatorname{DBD}(2n+1,3n+2)=1
    \]
    su kiekvienu sveikuoju $n$.

\end{enumerate}

\subsection{D lygis. Olimpiadinis taikymas}

\begin{enumerate}
    \renewcommand{\labelenumi}{\textbf{\theenumi.}}
    \setcounter{enumi}{15}

    \item Raskite visus sveikuosius skaičius $n$, kuriems
    \[
    n+5\mid n^2+3n+7.
    \]

    \item Raskite visus sveikuosius skaičius $n$, kuriems trupmenos
    \[
    \frac{n^2+4n+7}{n+2}
    \]
    reikšmė yra sveikasis skaičius.

    \item Įrodykite, kad bet kuriam sveikajam $n$
    \[
    \operatorname{DBD}(3n+4,5n+7)
    \]
    dalija skaičių $1$. Kokią išvadą galima padaryti apie šį DBD?

    \item Įrodykite: jeigu $a$ ir $b$ yra tarpusavyje pirminiai ir abu dalija
    sveikąjį skaičių $c$, tai
    \[
    ab\mid c.
    \]

    \item Tegu $a$ ir $b$ yra nelyginiai sveikieji skaičiai. Įrodykite, kad
    skaičius
    \[
    a^2-b^2
    \]
    dalijasi iš $8$.

\end{enumerate}

% ============================================================
\section*{Po treniruotės}
% ============================================================

Atsakykite į klausimus.

\begin{enumerate}
    \renewcommand{\labelenumi}{\textbf{\theenumi.}}

    \item Ką tiksliai reiškia žymėjimas $b\mid a$?

    \item Kodėl iš $d\mid a$ ir $d\mid b$ galime daryti išvadą, kad
    \[
    d\mid(ma+nb)
    \]
    visiems sveikiesiems $m,n$?

    \item Kodėl DBD uždaviniuose dažnai naudinga atimti skaičius?

    \item Ką pirmiausia verta tikrinti, kai uždavinyje klausiama, ar tam tikra
    situacija su sveikaisiais skaičiais yra įmanoma?

    \item Kurio uždavinio pagrindinė idėja šiandien buvo netikėčiausia?
\end{enumerate}

\end{document}
